\paragraph{Measured outcome.}
The warmed retained-LU solve adds zero allocations in the seven-size contract
cohort. The adapter removes 17 of 2119 and 50 of 2655 charged work units in the
8- and 16-dimensional fixtures, respectively, with exactly equal accepted
trajectories. All three dense fixtures likewise preserve every common observation
value and native counter. A separate Callgrind comparison measures instruction
reductions of 13.196\% for the 32-dimensional LU kernel, 2.605\% for the full
65-dimensional dense trajectory, and 1.731\% for the 16-dimensional adapter.
These establish implementation-work reductions in the profiled cohorts.
Wall-clock acceleration is \emph{not established}:
after an initial short-batch pass, a single longer confirmation gives paired
dense baseline/candidate ratios 1.033, 0.923 at dimensions 65 and 192 and
an adapter ratio 0.980 at dimension 16, but every observed block range
crosses unity. Process CPU ratios are also variable. Ratios above one favor
the candidate; neither attractive medians nor lower work counts are promoted
to reliable time-speedup claims. The shifted-factor research candidate is
excluded from production timing claims.
